\documentclass[conference]{IEEEtran}

% ─── Packages ───────────────────────────────────────────────────────
\usepackage[utf8]{inputenc}
\usepackage[T1]{fontenc}
\usepackage{amsmath, amssymb, amsfonts}
\usepackage{graphicx}
\usepackage{booktabs}
\usepackage{cite}
\usepackage{hyperref}
\usepackage{caption}
\usepackage{subcaption}
\usepackage{float}
\usepackage{enumitem}
\usepackage{multirow}
\usepackage{array}
\usepackage[ruled,vlined]{algorithm2e}
\usepackage{xcolor}

\hypersetup{colorlinks=true, linkcolor=blue, citecolor=blue, urlcolor=blue}

% ─── Title ──────────────────────────────────────────────────────────
\title{Parkinson's Disease Detection from Bilateral Wrist-Worn Smartwatch IMU Data Using Dual-Channel Cross-Attention Transformers}

\author{
  \IEEEauthorblockN{Author Name}
  \IEEEauthorblockA{\\
  University Name \\
  \texttt{author@email.edu}}
}

% ════════════════════════════════════════════════════════════════════
\begin{document}
\maketitle

% ════════════════════════════════════════════════════════════════════
% ABSTRACT
%   [✓] Problem statement
%   [✓] Method summary
%   [✓] Dataset description (PADS, 469 participants, 3 groups)
%   [✓] Key results (92.83%, SSL 20% labels, 48.32 ms edge)
%   [TODO] Significance/impact sentence — add 1 sentence at the end:
%          why does this matter clinically? e.g., "Early and continuous
%          monitoring of PD motor symptoms could enable timely intervention
%          and reduce reliance on specialist clinic visits."
%   [TODO] Keywords — add \begin{IEEEkeywords}...\end{IEEEkeywords}
%          after \end{abstract}. Suggested keywords:
%          Parkinson's Disease, IMU, wearable sensing, cross-attention
%          transformer, self-supervised learning, contrastive learning,
%          nested cross-validation, edge deployment
% ════════════════════════════════════════════════════════════════════
\begin{abstract}
% Parkinson's Disease (PD) is a chronic neurodegenerative disorder characterised by motor symptoms such as tremor, bradykinesia, postural instability, and freezing of gait (FoG). Clinical diagnosis relies on evaluation by a physician or specialist, which is time-consuming, subject to inter-rater variability, and inaccessible in many settings. Advanced wearable sensors offer a promising alternative by enabling continuous, objective monitoring of motor function in daily life. In this work, we propose a Dual-Channel Cross-Attention Transformer that processes bilateral wrist-worn IMU signals from the PADS dataset (469 participants: 291 PD, 79 HC, 99 DD) to classify Parkinson's Disease against both Healthy Controls (HC) and Differential Diagnosis (DD) conditions. Rather than a single three-class objective, we decompose the problem into two focused binary classification heads, enabling the model to learn the HC/PD and PD/DD decision boundaries independently. Patient-level stratified 5-fold cross-validation is applied throughout to obtain less biased estimates of true out-of-sample performance. The base model achieves a combined accuracy of $92.83\%$ across folds. We further demonstrate that self-supervised pre-training with contrastive infoNCE loss reaches $91.02\%$ combined accuracy with only $20\%$ of labeled data, highlighting the potential of label-efficient learning for clinical applications. The optimised model is successfully deployed on a Raspberry Pi, achieving a mean inference time of $48.32$~ms per window, confirming viability for real-time edge deployment on smartwatch hardware.
% [TODO] Add significance/impact sentence here before closing the abstract.
\end{abstract}

% [TODO] ADD KEYWORDS BLOCK HERE — required by IEEE:
% \begin{IEEEkeywords}
% Parkinson's disease, IMU, smartwatch, cross-attention transformer,
% self-supervised learning, infoNCE, nested cross-validation, edge deployment
% \end{IEEEkeywords}


% ════════════════════════════════════════════════════════════════════
% I. INTRODUCTION
%   [✓] Background (PD, symptoms, FoG, clinical diagnosis problem)
%   [✓] Problem statement (wearable sensors, difficulty of data mining)
%   [✓] Research gap (bilateral sensing neglected, label scarcity)
%   [✓] Contributions (4 bullet points)
%   [TODO] Paper organisation paragraph — MISSING. Every IEEE paper
%          must end the introduction with a paragraph like:
%          "The remainder of this paper is organised as follows.
%           Section~\ref{} describes the dataset. Section~\ref{}
%           details the methodology..." Add this as the final paragraph.
% ════════════════════════════════════════════════════════════════════
\section{Introduction}

%problem
Parkinson's Disease is a chronic neurodegenerative disorder affecting over 10 million people worldwide~\cite{ref}. It is characterised by progressive impairment of the dopaminergic system, leading to various motor symptoms such as resting tremor, bradykinesia, rigidity, postural instability, and freezing of gait (FoG). The symptoms duration can persist for a few seconds to several minutes, significantly increasing fall risk and reducing quality of life. Currently, PD is measured and analysed by a physician or specialist through which can take long hours.
%challenge
wearable sensors, specially smartwatch devices are embedded with sensors which can record IMU signals of motor activities. However, mining meaningful patterns and trends from such medical time series data is inherently challenging. Moreover, because of underlying similarity in motor patterns, distinguishing Parkinson's Disease(PD) from Differential Diagnosis (DD)  that produce Parkinson's-like symptoms is very difficult. 

%approaches so far
Deep learning approaches have shown strong performance in Human Activity Recognition (HAR) and medical time series classification. 
% Sequential models such as LSTMs capture temporal dependencies in sensor signals, while Transformer architectures with self-attention mechanisms are well suited to modelling long-range temporal patterns. In parallel, self-supervised contrastive representation learning has emerged as a powerful paradigm for learning from unlabeled data, which is especially relevant in clinical settings where obtaining true labels requires expert annotation and is therefore costly.

In this paper, we propose a Dual-Channel Cross-Attention Transformer framework for PD detection from bilateral wrist-worn IMU data.
\begin{itemize}
  % \item A dual-stream cross-attention architecture that explicitly models the motor asymmetry between the left and right wrists, which is a clinically established characteristic of Parkinson's Disease.
  % \item A hierarchical binary classification strategy that decomposes the three-group problem into two focused tasks (HC-PD and PD-DD), yielding a $4.48$~pp improvement in combined accuracy over a three-class baseline.
  % \item A self-supervised pre-training stage using infoNCE contrastive loss that achieves label-efficient transfer, reaching competitive performance with as few as $20\%$ of available labels.
  % \item Deployment of the optimised model on a Raspberry Pi CPU, demonstrating real-time inference at $48.32$~ms per window and validating feasibility for edge deployment on smartwatch hardware.
\end{itemize}

% [TODO] ADD PAPER ORGANISATION PARAGRAPH HERE. Example:
% The remainder of this paper is organised as follows.
% Section~\ref{} describes the PADS dataset and sensor modalities.
% Section~\ref{} details the proposed methodology, including preprocessing,
% model architecture, and training procedure. Section~\ref{} presents
% experimental results and ablation studies. Section~\ref{} discusses
% findings and limitations. Section~\ref{} concludes the paper.


% ════════════════════════════════════════════════════════════════════
% II. RELATED WORK
%   [✓] IMU-Based PD Detection subsection
%   [✓] Transformer Architectures for Medical Time Series subsection
%   [✓] Self-Supervised Contrastive Learning subsection
%   [✓] Nested Cross-Validation subsection
%   [TODO] All \cite{ref} placeholders must be replaced with real citations.
%          See supervisor notes for specific papers to cite per subsection.
%   [TODO] Each subsection should end with a sentence explicitly stating
%          the gap that YOUR work fills. Currently only the first subsection
%          does this clearly. Add gap sentences to subsections 2, 3, 4.
% ════════════════════════════════════════════════════════════════════
\section{Related Work}

% \subsection{IMU-Based Parkinson's Disease Detection}
% % [✓] Gap statement present at end of this subsection
% Wearable IMU sensors have been widely explored for objective PD assessment. Early work focused on handcrafted features extracted from accelerometer signals to detect tremor and bradykinesia~\cite{ref}. More recently, deep learning approaches have demonstrated competitive performance without manual feature engineering. Convolutional Neural Networks (CNNs) have been applied to IMU time series for gait and tremor classification, and recurrent architectures such as LSTMs have been used to capture the sequential dependencies in motor signals. Systematic reviews of deep learning with wearable sensors for PD~\cite{ref} cover datasets ranging from a handful of participants to over two thousand, highlighting the diversity of experimental protocols and the difficulty of cross-study comparison. A key limitation across this body of work is the predominant use of unilateral sensing, despite the well-established motor asymmetry of PD, where symptoms onset and progress more severely on one side of the body. Our work explicitly exploits bilateral sensing through a cross-attention mechanism designed around this clinical observation.

% \subsection{Transformer Architectures for Medical Time Series}
% % [TODO] Add gap sentence at the end: what does prior work NOT do that you do?
% %        e.g., "However, existing dual-stream designs fuse streams through
% %        concatenation or addition, without explicitly modelling the
% %        contralateral relationship between left and right motor signals."
% Transformer models, originally proposed by Vaswani et al.~\cite{vaswani2017} for natural language processing, have been increasingly applied to time series classification. Their self-attention mechanism allows the model to attend to distant timesteps simultaneously, which is advantageous for capturing the irregular, episodic nature of motor symptoms such as FoG. Dual-stream or two-stream transformer designs have been proposed for Human Activity Recognition (HAR), using parallel encoders to process different sensor modalities or spatial streams before fusion~\cite{ref}. Our architecture extends this paradigm to bilateral wrist sensing, replacing independent fusion with cross-attention between the two streams to explicitly model inter-wrist asymmetry.

% \subsection{Self-Supervised Contrastive Learning for Time Series}
% % [TODO] Add gap sentence at the end: e.g., "However, these methods have
% %        not been evaluated on bilateral IMU data for PD detection, where
% %        the augmentation strategy must preserve clinically meaningful
% %        asymmetric motor patterns across both wrist streams."
% Self-supervised contrastive representation learning has gained significant traction as a method for learning from unlabeled data. In the contrastive framework, a model is trained to map an anchor sample $x_i$ and its augmented positive counterpart $x_i^+$ to nearby embeddings, while pushing embeddings of negative samples $x_j$ ($j \neq i$) further apart. The infoNCE (Information Noise-Contrastive Estimation) loss~\cite{ref} formalises this objective by treating all other samples in a batch as negatives, enabling the model to contrast several negative samples simultaneously and making it more efficient for representation learning than triplet-based objectives.

% TS-TCC~\cite{ref} proposes a time-series-specific framework combining temporal and contextual contrasting with strong and weak augmentations, demonstrating competitive label efficiency on HAR and medical datasets. These methods show that self-supervised contrastive representation learning can learn general representations from unlabeled medical time series data that transfer effectively to downstream tasks with limited annotations. Our work applies this approach specifically to bilateral IMU data for PD detection, using temporal warping as the primary augmentation motivated by the non-linear temporal distortions characteristic of motor impairment.

% \subsection{Nested Cross-Validation in Clinical ML}
% % [TODO] Add gap sentence at the end: e.g., "Despite its importance,
% %        many published PD detection studies report performance from a
% %        single train/test split or flat k-fold, making their results
% %        difficult to compare and potentially overoptimistic."
% A persistent problem in medical machine learning is optimistic bias: when hyperparameter selection and performance estimation are performed on the same data split, reported metrics overestimate true out-of-sample performance. Nested cross-validation~\cite{ref} addresses this by separating the model selection loop (inner folds) from the performance estimation loop (outer folds), providing less biased estimates of how well predictive performance generalises to new data points. This bias-variance tradeoff is especially important in clinical settings with small patient cohorts, where single-fold or holdout validation can produce highly variable and misleading estimates. In this work, we apply Optuna-based hyperparameter search on the inner fold (Fold~1) and evaluate the selected configuration across all five outer folds using patient-level stratified splitting to prevent data leakage.


% ════════════════════════════════════════════════════════════════════
% III. DATASET
%   [✓] Dataset overview (PADS, PhysioNet, 469 participants, 3 groups)
%   [✓] Sensor modality (6-axis IMU, 100 Hz)
%   [✓] Motor tasks (10 tasks listed)
%   [TODO] Add a sentence on task duration — Relaxed is 20s, all others
%          10s. This is relevant to your windowing strategy in methodology.
%   [TODO] Replace \cite{ref} with actual PADS dataset citation.
% ════════════════════════════════════════════════════════════════════
\section{Dataset}

\subsection{Overview}
For this experiment, we have used the PADS (Parkinson's Disease Smartwatch) dataset from PhysioNet~\cite{ref}. It contains sensor data from bilateral wrist-worn smartwatch recordings of 469 participants performing $10$ standardized motor tasks. The dataset comprises mainly three groups: 291 patients with Parkinson's Disease (PD), 79 Healthy Controls (HC), and 99 patients with Differential Diagnosis (DD) conditions that show Parkinson's-like motor symptoms but are different neurodegenerative diseases. Inertial Measurement Unit (IMU) data is highly effective in detecting motor symptoms such as tremor and bradykinesia and highly applicable for this task.

\subsubsection{Sensor Modality}
Each participant wears a smartwatch on both wrists, providing six-axis IMU measurements per wrist. Three axes are from the accelerometer ($a_x, a_y, a_z$), which detects linear movements, and three axes are from the gyroscope ($\omega_x, \omega_y, \omega_z$), which records angular movements. The signals are sampled at $100$~Hz frequency.
Figure~\ref{fig:sensor_viz} visualizes the sensor data of a patient performing the \textit{HoldWeight} task, showing distinct patterns in the accelerometer and gyroscope channels indicative of motor symptoms such as tremor and bradykinesia.

\begin{figure}[H]
\centering
\includegraphics[width=0.45\textwidth]{image5.png}
\caption{IMU sensor data from a participant performing the \textit{HoldWeight} task.}
\label{fig:sensor_viz}
\end{figure}

\subsubsection{Motor Tasks}
% [TODO] Add task duration info: "Tasks \textit{Relaxed} and
%        \textit{RelaxedTask} are 20~seconds in duration; all remaining
%        tasks are 10~seconds." This affects how many windows each task yields.
Ten standardized motor tasks are recorded from clinical trials: \textit{CrossArms}, \textit{DrinkGlas}, \textit{Entrainment}, \textit{HoldWeight}, \textit{LiftHold}, \textit{PointFinger}, \textit{Relaxed}, \textit{StretchHold}, \textit{TouchIndex}, and \textit{TouchNose}. 


% ════════════════════════════════════════════════════════════════════
% IV. METHODOLOGY
%   [✓] Problem Definition (corrected notation X ∈ R^{T×2C})
%   [✓] Data Preprocessing
%       [✓] Label Encoding
%       [✓] K-Fold Cross-Validation (corrected: reduces variance, not increases)
%       [✓] Downsampling
%       [✓] Band-Pass Filtering
%       [✓] Window Overlapping (with formula)
%   [✓] Model Architecture
%       [✓] Input Projection and Positional Encoding
%       [✓] Cross-Attention Layers (corrected Q/K/V description + equation)
%       [✓] Pooling, Fusion, Classification
%   [✓] Training Procedure
%       [✓] Objective Function
%       [✓] Optimization
%   [✓] Self-Supervised Pre-training
%       [✓] Framework description (anchor/positive/negative, infoNCE)
%       [✓] Encoder-only pre-training stated explicitly
%       [✓] Augmentation (time warping + Gaussian noise)
%       [✓] infoNCE loss equation
%   [TODO] Add citation for time-warping augmentation (Um et al. 2017)
%   [TODO] Add citation for infoNCE loss (Oord et al. 2018)
%   [TODO] Add citation for AdamW optimizer
% ════════════════════════════════════════════════════════════════════
\section{Methodology}

\subsection{Problem Definition}

This study aims to develop a method for classifying Parkinson's Disease (PD) from Healthy Controls (HC) and Differential Diagnosis (DD) using time-series data collected from wearable IMU sensors.
Each input sample contains left ($L$) and right ($R$) wrist signal pairs, which we can denote as:
$\mathbf{X} = (\mathbf{X}^L, \mathbf{X}^R) \in \mathbb{R}^{T \times 2C}$,
where $C{=}6$ denotes the number of sensor channels per wrist, $2C{=}12$ is the total number of channels across both wrists, and $T{=}256$ is the total number of timesteps in a sample.
Each sample is associated with a class label $y \in \{\text{HC}, \text{PD}, \text{DD}\}$.

Instead of considering the samples as three independent classes, we divide the classification problem into two distinct hierarchical binary classifiers. The model learns a mapping function
$f : \mathbb{R}^{T \times 2C} \rightarrow [0,1]^2$, where the model predicts two probabilities corresponding to HC-PD (from the sample space $y \in \{\text{HC}, \text{PD}\}$) and PD-DD (from the sample space $y \in \{\text{PD}, \text{DD}\}$).
The PD-DD boundary shares similar motor patterns, Making partitioning this second decision boundary more challenging.

\subsection{Data Preprocessing}

\subsubsection{Label Encoding}
A hierarchical labeling strategy is applied to support two independent binary classification heads. DD samples are masked ($-1$) and ignored during loss computation for the HC-PD head; HC samples are masked ($-1$) for the PD-DD classification. Each classification head is trained on relevant samples and ignores samples unrelated to that class which prevents data leakage.

\subsubsection{K-Fold Cross-Validation}
A \textbf{patient-level stratified $5$-fold cross-validation} strategy is applied. The K-fold splits the dataset into 5 equal parts, and the stratified cross-validation ensures equal distribution of samples from different classes in each fold. K-fold cross-validation is applied to obtain less biased estimates of the true out-of-sample predictions. A patient-level train/test split is applied: all windows from a given patient are assigned entirely to either the training or the validation set, preventing data leakage across patient boundaries.

\subsubsection{Downsampling}
Raw signals are resampled from $100$~Hz to $64$~Hz via polyphase resampling. Because higher frequency signal is irrelevant to HAR activities, as human motor signals are most of the time below $20$~Hz.

\subsubsection{Band-Pass Filtering}
Human voluntary movements and clinically relevant motor symptoms such as tremor and FoG are concentrated in the low-frequency, high-sensitivity band of $0.1$--$20$~Hz. A fourth-order Butterworth band-pass filter with passband $[0.1,\,20]$~Hz is applied using zero-phase forward-backward filtering to avoid phase distortion. The low-frequency cut-off removes DC drift, and the high-frequency cut-off suppresses sensor noise above the motor-relevant band.

\subsubsection{Window Overlapping}
\label{sec:windowing}
We have segmented the IMU signals into equal length consecutive windows. Each window sample is $256$ timestamps, or $4$~seconds, long (at $64$~Hz). As the number of patients in the trial is disproportionate (HC: 79, PD: 291, DD: 99); applying uniform non-overlapping window sampling produces severe class imbalance, resulting in the model biasing toward the PD majority class. To mitigate this, an \textbf{overlapping sampling} strategy is applied: minority classes receive higher window overlap, generating more samples from the same recording duration. We applied $70$\% overlap to HC samples, $0$\% overlap to PD, and $65$\% overlap to DD. This resulted in approximately $8{,}000$ windows per class in the training set ($25{,}640$ total), distributed evenly across folds via stratified cross-validation. If the number of windows produced for a signal of length $L$ with window size $W$ and step size is $s$  we can represent it as:
\begin{equation}
  N_w = \left\lfloor \frac{L - W}{s} \right\rfloor + 1
\end{equation}

\subsection{Dual-Channel Cross-Attention Transformer}
\label{sec:base_model}
The architecture processes left and right wrist window samples independently through parallel encoders. The design is motivated by the well-established motor asymmetry of Parkinson's Disease: symptoms typically occur more severely on one side, making cross-attention between the left and right wrist signals clinically meaningful. Fig.~\ref{fig:architecture} illustrates the full pipeline.

\begin{figure*}[t]
\centering
\includegraphics[width=\textwidth]{cross_attn.png}
\caption{Dual-channel cross-attention Transformer.}
\label{fig:architecture}
\end{figure*}

\subsubsection{Input Projection and Positional Encoding}
Each wrist's $C{=}6$ channel input is linearly projected to a $d$-dimensional embedding space. Sinusoidal positional encodings are applied to inject temporal order into the embeddings, following the standard transformer methodology~\cite{vaswani2017}.

\subsubsection{Cross-Attention Layers}
Each of the cross-attention module has three sequential sub-layers. First, in the cross-attention sub-layer, each wrist's embedding acts as Queries ($Q$) and attends over the opposite wrist's embedding as Keys ($K$) and Values ($V$) and vice versa. The output embedding is followed by a residual connection and layer normalization. For the left stream, where $\mathbf{H}^L$ queries into $\mathbf{H}^R$:
\begin{equation}
  \mathbf{H}^L \leftarrow \mathrm{LayerNorm}\!\left(
      \mathbf{H}^L + \mathrm{MHA}(\mathbf{H}^L,\, \mathbf{H}^R,\, \mathbf{H}^R)
  \right)
\end{equation}
The right stream is updated symmetrically with $\mathbf{H}^R$ as queries over $\mathbf{H}^L$. Second, a \emph{self-attention} sub-layer is applied to each wrist's own embedding. Third, a position-wise \emph{feed-forward} network containing a two-layer MLP with ReLU activation is applied.

\subsubsection{Pooling, Fusion, and Classification}
After $N$ attention layers, global average pooling is applied to each stream, yielding $\mathbf{z}^L, \mathbf{z}^R \in \mathbb{R}^d$. The two embeddings are concatenated to form $\mathbf{z} = [\mathbf{z}^L \| \mathbf{z}^R] \in \mathbb{R}^{2d}$, which is fed into two independent classification heads --- one for HC-PD and one for PD-DD --- each implemented as a single linear layer followed by softmax.

\subsection{Training Procedure}

\subsubsection{Objective Function}
The cross-entropy loss is applied in each classification head, computed only over samples with valid (non-masked) labels of each classifier. The joint loss is the average of the two masked losses:
\begin{equation}
\mathcal{L} = \frac{1}{|\mathcal{S}|} \sum_{s \in \mathcal{S}} \mathcal{L}_s, \qquad
\mathcal{L}_s = -\sum_{c=1}^{2} y_{s,c} \log \hat{p}_{s,c}
\end{equation}

\subsubsection{Optimization}
The model is trained with AdamW (weight decay $0.01$), gradient clipping at unit norm, and a ReduceLROnPlateau scheduler (factor $0.5$, patience $5$ epochs). Training runs for up to $100$ epochs with batch size $32$.

\subsection{Self-Supervised Pre-training}
\label{sec:ssl}
We applied self-supervised contrastive representation learning as a pre-training stage before fine-tuning on labelled data. True labels for clinical samples are scarce. By first learning general signal representations from unlabelled data, the model can then transfer useful structural information to downstream classification task.

Three types of samples are required in contrastive learning; anchor sample  $z_i$,positive sample $z_i^+$ created by augmenting the anchor  and negative sample $z_j$ ($j \neq i$).  
We have applied two augmentation strategy; time warping and gaussian noise augmentation alternating with a $50$\% probability. with Random Sampling Strategy and contrastive \textbf{infoec loss}. 
the model attempts to find similarity between anchor and positive samples and push negative samples.
The time warping augmentation is perturbed via cubic spline interpolation with $k{=}4$ interior knots displaced from $\mathcal{N}(1.0,\, 0.3^2)$. Same augmentations was applied to all three types of samples and the gaussian noise with  $strength = 0.05$ was applied.
The infoNCE(Information Noise-Contrastive Estimation ) loss compels embedding of an anchor sample   $z_i$  to be closer to its corresponding positive sample   $z_i^+$ while pushing it away from other samples in the batch. infoNEC loss can compare several negative samples with the anchor samples which is more efficient in case of learning representation.
\begin{equation}
\mathcal{L} = 
- \frac{1}{B} \sum_{i=1}^{B}
\log \frac{
\exp\left(\mathrm{sim}(z_i, z_i^+)/\tau\right)
}{
\sum_{j=1}^{B} \exp\left(\mathrm{sim}(z_i, z_j)/\tau\right)
}
\end{equation}

where $\mathrm{sim}(\cdot,\cdot)$ denotes cosine similarity and $\tau$ is a temperature hyperparameter. After pre-training, the encoder is fine-tuned end-to-end on the labeled downstream task.


% ════════════════════════════════════════════════════════════════════
% V. EXPERIMENTS AND RESULTS
%   [✓] Experimental setup (described in Implementation Details section)
%   [✓] Evaluation metrics (Acc, Prec, Rec, F1, AUC via ROC curves)
%   [✓] Base model performance (5-fold mean±std reported)
%   [✓] Band-pass filtering ablation
%   [✓] Hierarchical vs three-class comparison
%   [✓] Hyperparameter optimisation (Optuna, nested CV)
%   [✓] Edge deployment (Raspberry Pi, 48.32 ms)
%   [✓] SSL label efficiency (Table + Figure)
%   [✓] TimesFM LoRA comparison
%   [✓] Task-wise ablation (CNN vs BiLSTM)
%   [✓] Attention weight analysis
%   [TODO] MISSING: Comparison with state-of-the-art on PADS dataset.
%          Add a table comparing your method against prior published results
%          on the same dataset (or closest available). Without this,
%          reviewers cannot assess whether 92.83% is a strong result.
%   [TODO] MISSING: Per-fold breakdown table. Add a table like:
%          Fold | HC-PD Acc | PD-DD Acc | Combined
%            1  |  ...      |  ...      |  ...
%            2  |  ...      |  ...      |  ...
%          so readers can verify consistency across folds, not just mean±std.
%   [TODO] Setup subsection is scattered across Implementation Details
%          (at the end) and inline text. Consider adding a short
%          \subsection{Experimental Setup} at the start of this section.
% ════════════════════════════════════════════════════════════════════
\section{Experiments and Results}

% [TODO] ADD: \subsection{Experimental Setup}
% Move/summarise from Implementation Details: hardware, framework,
% evaluation metrics definition, and CV strategy used for all experiments.

\subsection{Base Model Performance}
The base model trained with band-pass filtering and downsampling achieved a mean accuracy of $93.31\%\pm0.22\%$ (precision $93.81\%\pm0.31\%$, recall $93.31\%\pm0.22\%$, F1 $93.23\%\pm0.22\%$) for HC-PD classification and $92.34\%\pm0.08\%$ accuracy (precision $93.14\%\pm0.18\%$, recall $92.34\%\pm0.08\%$, F1 $92.28\%\pm0.08\%$) for PD-DD classification across 5-fold cross-validation, giving a combined accuracy of $92.83\%$. The model showed higher accuracy for the HC-PD partition than the PD-DD partition, which confirms that differentiating Parkinson's patients (PD) from differential diagnosis (DD) is more challenging due to the commonality of motor patterns. From the t-SNE plot Fig.~\ref{fig:initial_tsne} we can observe that there exists a particular region common for both PD-DD samples, unlike HC-PD samples.

\begin{figure}[H]
\centering
\begin{subfigure}[b]{0.23\textwidth}
    \centering
    \includegraphics[width=\textwidth]{roc_hc_vs_pd.png}
    \caption{HC vs PD}
\end{subfigure}
\hfill
\begin{subfigure}[b]{0.23\textwidth}
    \centering
    \includegraphics[width=\textwidth]{roc_pd_vs_dd.png}
    \caption{PD vs DD}
\end{subfigure}
\caption{ROC curves, base model (Fold~1).}
\label{fig:initial_roc}
\end{figure}

\begin{figure}[H]
\centering
\begin{subfigure}[b]{0.23\textwidth}
    \centering
    \includegraphics[width=\textwidth]{tsne_hc_vs_pd.png}
    \caption{HC vs PD}
\end{subfigure}
\hfill
\begin{subfigure}[b]{0.23\textwidth}
    \centering
    \includegraphics[width=\textwidth]{tsne_pd_vs_dd.png}
    \caption{PD vs DD}
\end{subfigure}
\caption{t-SNE embeddings, base model (Fold~1).}
\label{fig:initial_tsne}
\end{figure}

% [TODO] ADD PER-FOLD BREAKDOWN TABLE HERE. Example structure:
% \begin{table}[H]
% \centering
% \caption{Per-fold results for the base model.}
% \begin{tabular}{ccccc}
% \toprule
% \textbf{Fold} & \textbf{HC-PD Acc} & \textbf{PD-DD Acc} & \textbf{Combined} \\
% \midrule
% 1 & ? & ? & ? \\
% 2 & ? & ? & ? \\
% 3 & ? & ? & ? \\
% 4 & ? & ? & ? \\
% 5 & ? & ? & ? \\
% \midrule
% Mean$\pm$Std & $93.31\pm0.22$ & $92.34\pm0.08$ & $92.83\%$ \\
% \bottomrule
% \end{tabular}
% \end{table}

\subsection{Band-Pass Filtering}
\label{sec:ablation_bp}
Band-pass filtering is fundamental in physiological signal processing. Significant human motor activity such as tremor and freezing of gait is concentrated in the $0.1$--$20$~Hz band; signals with Frequency above this limit are primarily from sensor noise unrelated to motor symptoms.
Without band-pass filtering, the model achieved a mean accuracy of only $85.64\%\pm2.14\%$ (HC-PD) and $85.39\%\pm1.88\%$ (PD-DD) with very high variance across folds, and the training curves in Fig.~\ref{fig:bandpass_loss} confirm severe over fitting. With the filter applied, accuracy improved to $93.31\%\pm0.22\%$ and $92.34\%\pm0.08\%$ respectively, with significantly reduced fold-to-fold variance.
% demonstrating that removing out-of-band noise is essential for reliable learning.

\begin{figure}[H]
\centering
\begin{subfigure}[b]{0.23\textwidth}
    \centering
    \includegraphics[width=\textwidth]{loss_no_bandpass.png}
    \caption{Without band-pass}
\end{subfigure}
\hfill
\begin{subfigure}[b]{0.23\textwidth}
    \centering
    \includegraphics[width=\textwidth]{loss.png}
    \caption{With band-pass}
\end{subfigure}
\caption{Loss curves illustrating overfitting without band-pass filtering (Fold~1).}
\label{fig:bandpass_loss}
\end{figure}

\subsection{Hierarchical vs.\ Three-Class Classification}
The three-class model achieved a mean accuracy of $88.35\%$. The learning objective of  distinguishing healthy samples(HC) from diseased samples(PD, DD) is different than distinguishing  between PD and DD. The model fails to focus on the subtle difference between PD and DD samples resulting in lower accuracy in PD and DD class. Fig.~\ref{fig:3class_eval} we can observe that the model struggles to separate PD and DD samples relative to HC.

\begin{figure}[H]
\centering
\begin{subfigure}[b]{0.23\textwidth}
    \centering
    \includegraphics[width=\textwidth]{roc_3class.png}
    \caption{Three-class ROC}
\end{subfigure}
\hfill
\begin{subfigure}[b]{0.23\textwidth}
    \centering
    \includegraphics[width=\textwidth]{tsne_3class.png}
    \caption{Three-class t-SNE}
\end{subfigure}
\caption{Three-class baseline. The t-SNE shows significant PD/DD overlap.}
\label{fig:3class_eval}
\end{figure}

\subsection{Hyperparameter Optimisation}

Hyperparameter tuning was done to achieve regularization and to improve generalization to unseen data. Nested cross-validation is applied to reduce optimistic bias of the model.
Hyperparameter search was conducted on Fold~1 as the internal validation fold within the nested cross-validation set; the remaining folds serve as holdout validation to estimate generalisation performance. This ensures that the model is a less biased estimate of true out-of-sample performance.

\begin{table}[H]
\centering
\caption{Training vs.\ Testing Accuracy at Key Epochs (Fold~1, Optuna Model)}
\label{tab:train_test_gap}
\begin{tabular}{lccc}
\toprule
\textbf{Phase} & \textbf{Train} & \textbf{Test} & \textbf{Gap} \\
\midrule
 (Epoch 45)& 89.39\% & 91.82\% & \textbf{$-$2.43\%} \\
 (Epoch 80)& 92.10\% & 88.78\% & $+$3.32\% \\
\bottomrule
\end{tabular}
\end{table}

The optimised model (Epoch~45) achieves a test accuracy of $88.78\%$, which is approximately $2.5$~pp below the base model, in exchange for a $2\times$ reduction in embedding dimension ($d{=}32$ vs.\ $d{=}64$). The negative gap at Epoch~45 (test accuracy exceeds training accuracy) indicates good generalization at that point rather than overfitting. The optimised model also shows comparatively stronger performance on the PD-DD classification relative to the base model, which proves that smaller,  regularised model has high variance to distinnguish subtle features from PD/DD class./*

\subsection{Edge Deployment on Raspberry Pi}
\label{sec:pi}
The model was deployed on Raspberry Pi CPU to observe real-time inference. The inference pipeline downsamples incoming $100$~Hz samples to $64$~Hz, applies a causal fifth-order Butterworth band-pass filter sample-by-sample using second-order sections (SOS) with persistent filter state, and accumulates $256$ samples in a circular buffer. Inference is triggered every $128$ new samples ($50$\% overlap, $2.0$~s step), producing a prediction at $0.5$~Hz.
The inference showed an average per-window inference time of $48.32$~ms with a minimum CPU utilisation of $0.02\%$. This observation provides credibility for deployment on an edge device such as a smartwatch for real-time detection.

\subsection{Self-Supervised Pre-training }
The self-supervised model fine-tuned on labelled data achieved a combined accuracy of $92.87\%$ (HC-PD: $93.35\%$, PD-DD: $92.40\%$). A label-efficiency study evaluated fine-tuning across fractions of the labelled set; results are summarised in Table~\ref{tab:ssl_label_eff}. With only $5\%$ of labelled data the model reaches $80.49\%$ combined accuracy, and with $20\%$ it already exceeds $91\%$ and a much slower gain from $20\%$ to $100\%$ ($91.02\%$ to $92.87\%$, a gain of less than $2$~pp), which demonstrates that the self-supervised pre-training has learned general representations and  requires very few labels to reach saturation.
% transfer effectively to the downstream classification task.

\begin{table}[H]
\centering
\caption{SSL label-efficiency: fine-tuning accuracy vs.~labelled fraction.}
\label{tab:ssl_label_eff}
\begin{tabular}{rcccc}
\toprule
\textbf{Labels (\%)} & \textbf{HC vs PD} & \textbf{PD vs DD} & \textbf{Combined} \\
\midrule
$5\%$   & $80.52\%$ & $80.46\%$ & $80.49\%$ \\
$10\%$  & $87.62\%$ & $86.12\%$ & $86.87\%$ \\
$20\%$  & $92.22\%$ & $89.82\%$ & $91.02\%$ \\
$50\%$  & $92.70\%$ & $91.69\%$ & $92.20\%$ \\
$70\%$  & $92.34\%$ & $92.10\%$ & $92.22\%$ \\
$100\%$ & $93.35\%$ & $92.40\%$ & $92.87\%$ \\
\bottomrule
\end{tabular}
\end{table}

\begin{figure}[H]
\centering
\includegraphics[width=0.45\textwidth]{accuracy_vs_labels.png}
\caption{Combined accuracy vs.\ labelled fraction.}
\label{fig:label_eff}
\end{figure}

\subsection{Finetuning TimesFM Model}
% [TODO] Add \cite for TimesFM and for LoRA (Hu et al. 2022)
For a comparative analysis we have fine-tuned the foundation model TimesFM~\cite{ref}, pre-trained on large datasets for general-purpose forecasting tasks. We adapted the model for classification by adding a LoRA (Low-Rank Adaptation) adapter~\cite{ref} with rank $r{=}6$, scaling factor $\alpha{=}12$, dropout $0.15$, and mixed regularization $0.1$, to the hidden dimension of $1280$. LoRA is a parameter-efficient fine-tuning technique that adapts large pre-trained models to new tasks on smaller datasets by learning low-rank update matrices rather than modifying all parameters. This training achieved $89.70\%$ accuracy for the HC-PD classifier and $87.80\%$ accuracy for the PD-DD classifier ($88.75\%$ combined), which is below the proposed base model, suggesting that a forecasting foundation model does not transfer as effectively to this classification setting as a task-specific architecture trained from scratch.

\subsection{Task-Wise Ablation (CNN and BiLSTM)}
We have trained task-specific 1D CNN and BiLSTM models for comparison. For the 1D CNN we used four convolutional blocks($d=128$) with kernel sizes $7, 5, 3, 3$ respectively, followed by batch normalization and max pooling. A task-wise bidirectional-LSTM model with two layers ($d=128$) followed by layer normalization was trained.

The BiLSTM showed higher accuracy for tasks such as CrossArms ($92.39\%$), HoldWeight ($93.87\%$), LiftHold ($93.92\%$), and StretchHold ($93.87\%$), which involves high temporal dependencies. The CNN showed comparatively performed higher on tasks such as Entrainment ($63.68\%$) and PointFinger ($65.86\%$), where spatial features within a window are more informative than temporal features. Tasks such as Entrainment, PointFinger, TouchIndex, and TouchNose consistently produced lower accuracy ($<70\%$) for both models, indicating that these tasks has motor patterns that are harder to discriminate at a window level for either architecture.

\begin{figure}[H]
\centering
\includegraphics[width=0.45\textwidth]{image6.png}
\caption{Task-wise ablation: 1D CNN vs.\ BiLSTM combined accuracy per task.}
\label{fig:taskwise}
\end{figure}

% [TODO] MISSING: Comparison with state-of-the-art table.
% Add a table here comparing your base model against published results
% on the PADS dataset or closely related PD+wearable+IMU papers.
% Format example:
% \begin{table}[H]
% \centering
% \caption{Comparison with prior work.}
% \begin{tabular}{lccc}
% \toprule
% Method & HC-PD Acc & PD-DD Acc & Combined \\
% \midrule
% Prior work A \cite{} & ? & ? & ? \\
% Prior work B \cite{} & ? & ? & ? \\
% \textbf{Ours} & 93.31 & 92.34 & \textbf{92.83} \\
% \bottomrule
% \end{tabular}
% \end{table}

\subsection{Attention Weight Analysis}
The cross-attention weight maps are visualized to provide understanding of which temporal patterns the model attends to. Figure~\ref{fig:attn} shows that attention bands become progressively sparser across the classes from HC to PD and PD to DD. The HC class shows concentrated activity on earlier timesteps and at timestep 250, while PD samples show more distributed active bands. This suggests that motor symptoms such as tremor and bradykinesia produce recurring temporal patterns that the model detects through sparse, distributed attention. The similarity between the active regions of PD and DD samples further confirms the difficulty of separating Parkinson's from other neurodegenerative disorders at the attention level.

\begin{figure}[H]
\centering
\begin{subfigure}[b]{0.45\textwidth}
    \centering
    \includegraphics[width=\textwidth]{attention_comparison_hc_vs_pd.png}
    \caption{HC vs PD}
\end{subfigure}
\vspace{4pt}
\begin{subfigure}[b]{0.45\textwidth}
    \centering
    \includegraphics[width=\textwidth]{attention_comparison_pd_vs_dd.png}
    \caption{PD vs DD}
\end{subfigure}
\caption{Cross-attention weight maps (Fold~1).}
\label{fig:attn}
\end{figure}


% ════════════════════════════════════════════════════════════════════
% VI. DISCUSSION
%   [✓] Interpretation of hierarchical classification result
%   [✓] Interpretation of band-pass filtering result
%   [✓] Interpretation of SSL label efficiency result
%   [✓] Clinical relevance of edge deployment
%   [✓] Task-wise ablation interpretation
%   [TODO] MISSING: Limitations subsection. This is the most commonly
%          missing section in student papers and the first thing reviewers
%          ask for. Add at minimum:
%          1. Dataset limitation: single dataset (PADS), no external validation
%          2. Generalisation: results may not hold for other smartwatch hardware
%          3. SSL pre-training: same dataset used for pre-training and fine-tuning
%             (ideally pre-training should be on separate unlabeled data)
%          4. Clinical validity: no clinician-in-the-loop validation of predictions
% ════════════════════════════════════════════════════════════════════
\section{Discussion}

The results demonstrate that the proposed Dual-Channel Cross-Attention Transformer is well suited to the interdisciplinary challenge of PD detection from wearable IMU data. The hierarchical binary classification strategy consistently outperforms the three-class baseline by decomposing a single hard objective into two focused tasks, which reflects the clinical reality that HC/PD and PD/DD represent qualitatively different discrimination problems. The $4.48$~pp combined accuracy gain confirms that the model benefits from learning each decision boundary independently.

The band-pass filtering ablation highlights a fundamental property of medical time series data: the discriminative signal is concentrated in a narrow, low-frequency, high-sensitivity band ($0.1$--$20$~Hz), and out-of-band noise does not merely add difficulty but actively prevents the model from learning meaningful patterns. The substantially higher fold-to-fold variance in the unfiltered condition ($\pm2.14$\% vs.\ $\pm0.22$\% for HC-PD) further illustrates that noise inflates variance in model estimates, which is precisely the regime where nested cross-validation and careful signal processing matter most.

The self-supervised pre-training results reveal an important property for clinical deployment: the model achieves over $91\%$ combined accuracy with just $20\%$ of labeled data, and further gains from additional labels are marginal. In practice, obtaining labeled clinical data requires a physician or specialist to evaluate each participant, which is a significant bottleneck. The SSL framework demonstrates that general IMU representations can be learned without any labels, substantially reducing the annotation burden for future data collection in new clinical populations or settings.

The edge deployment on Raspberry Pi validates the feasibility of real-time risk assessment on constrained hardware. At $48.32$~ms per inference window and $0.02\%$ CPU utilisation, the model is far within the computational budget of a modern smartwatch processor. This positions the framework as a viable component of a continuous disease monitoring system in which PD progression can be tracked passively during daily life, rather than relying on infrequent clinical visits.

The task-wise BiLSTM and CNN ablation reveals that no single baseline architecture is uniformly superior across all motor tasks, and that some tasks ($<70\%$ for both models) remain genuinely difficult regardless of architecture choice. This suggests that future work should investigate task-specific or task-aware models, and that the aggregation of multiple motor tasks --- rather than reliance on any single task --- is important for robust classification.

% [TODO] ADD LIMITATIONS SUBSECTION HERE:
% \subsection{Limitations}
% The primary limitation of this work is that all experiments are conducted
% on a single dataset (PADS). While patient-level cross-validation reduces
% optimistic bias within this dataset, external validation on an independent
% cohort is necessary to confirm generalisability. Additionally, the
% self-supervised pre-training uses the same dataset as fine-tuning;
% ideally, pre-training would leverage a larger pool of unlabeled IMU data
% from diverse populations. Finally, clinical validation --- where model
% predictions are evaluated by a specialist against ground-truth diagnosis
% --- has not been performed and remains essential before real-world deployment.


% ════════════════════════════════════════════════════════════════════
% VII. SUMMARY OF RESULTS
%   [✓] Consolidated table (mean±std across all variants)
%   [✓] Ablation comparison figure
%   [TODO] Consider merging this section into Experiments & Results
%          or Discussion. A standalone "Summary of Results" section is
%          unusual in IEEE conference papers and may use up page budget.
%          Reviewers may question why it is a separate section.
% ════════════════════════════════════════════════════════════════════
\section{Summary of Results}

Table~\ref{tab:summary} consolidates the 5-fold cross-validation results (mean$\,\pm\,$std\,\%) across all experimental variants. The hierarchical dual-head base model with band-pass filtering achieves the highest combined accuracy. Self-supervised pre-training matches the supervised baseline while requiring significantly fewer labels.

\begin{table*}[t]
\centering
\caption{Consolidated 5-fold cross-validation results (mean$\,\pm\,$std\,\%). Best combined accuracy in bold.}
\label{tab:summary}
\begin{tabular}{lccccccc}
\toprule
\multirow{2}{*}{\textbf{Approach}} &
  \multicolumn{3}{c}{\textbf{HC vs PD}} &
  \multicolumn{3}{c}{\textbf{PD vs DD}} &
  \multirow{2}{*}{\textbf{Combined}} \\
\cmidrule(lr){2-4}\cmidrule(lr){5-7}
& \textbf{Acc} & \textbf{Prec} & \textbf{Rec}
& \textbf{Acc} & \textbf{Prec} & \textbf{Rec} & \\
\midrule
Base Model (w/ BP)   & $93.31{\pm}0.22$ & $93.81{\pm}0.31$ & $93.31{\pm}0.22$ & $92.34{\pm}0.08$ & $93.14{\pm}0.18$ & $92.34{\pm}0.08$ & $\mathbf{92.83\%}$ \\
Base Model (w/o BP)  & $85.64{\pm}2.14$ & $85.71{\pm}2.15$ & $85.64{\pm}2.14$ & $85.39{\pm}1.88$ & $85.64{\pm}1.92$ & $85.39{\pm}1.88$ & $85.52\%$ \\
Three-Class Baseline & \multicolumn{3}{c}{Acc: $88.35{\pm}0.83$, Prec: $89.67{\pm}0.83$} & \multicolumn{3}{c}{Rec: $88.35{\pm}0.83$, F1: $88.32{\pm}0.88$} & $88.35\%$ \\
SSL Fine-tune (100\%)& $93.35$ & $93.96$ & $93.35$ & $92.40$ & $93.31$ & $92.40$ & $92.87\%$ \\
TimesFM LoRA ($r{=}6$) & $89.70$ & --- & --- & $87.80$ & --- & --- & $88.75\%$ \\
\midrule
Raspberry Pi (inference)  & \multicolumn{6}{c}{---} & $\approx48.32$~ms/win \\
\bottomrule
\end{tabular}
\end{table*}

\begin{figure}[H]
\centering
\includegraphics[width=0.45\textwidth]{plot_ablation_comparison.png}
\caption{Combined accuracy across all model variants.}
\label{fig:ablation}
\end{figure}


% ════════════════════════════════════════════════════════════════════
% VIII. CONCLUSION
%   [✓] Summary of contributions
%   [✓] Key results restated
%   [✓] Future work directions
%   [TODO] MISSING: Explicit limitations sentence in conclusion.
%          Add 1-2 sentences acknowledging the main limitation
%          (single-dataset, no external validation) before future work.
%   [TODO] MISSING: Closing clinical significance sentence.
%          End with why this matters for patients, not just for the model.
% ════════════════════════════════════════════════════════════════════
\section{Conclusion}

In this work, we proposed a Dual-Channel Cross-Attention Transformer framework for detecting Parkinson's Disease from bilateral wrist-worn IMU data. By processing left and right wrist signals through parallel encoders connected by cross-attention, the model explicitly captures the motor asymmetry that is a clinically established characteristic of PD. A hierarchical binary classification strategy decomposes the three-group problem into focused HC-PD and PD-DD tasks, yielding a $4.48$~pp improvement over a three-class baseline. Patient-level stratified 5-fold cross-validation throughout the experimental pipeline provides less biased estimates of true out-of-sample performance, and band-pass filtering is shown to be an essential preprocessing step, reducing fold-to-fold variance by an order of magnitude.

The self-supervised contrastive pre-training stage demonstrates that competitive classification performance is achievable with only $20\%$ of labeled data, which is a critical finding for developing clinical frameworks in which expert annotation is scarce. The model is successfully deployed on a Raspberry Pi at $48.32$~ms per inference window, confirming that real-time PD monitoring on edge wearable hardware is computationally feasible.

% [TODO] ADD LIMITATIONS SENTENCE HERE before future work. Example:
% "The primary limitation of this work is that evaluation is conducted
% on a single dataset; external validation on independent cohorts is
% required before clinical deployment."

To address the remaining challenges, future work should investigate task-aware fusion strategies to improve performance on difficult motor tasks such as Entrainment and PointFinger where both CNN and BiLSTM models underperform. Expanding the framework to include additional sensor modalities and longitudinal data would enable risk assessment and disease progression monitoring over time. The interdisciplinary nature of this problem --- spanning clinical neurology, signal processing, and deep learning --- means that clinical validation in partnership with specialist assessment will be an essential next step before real-world deployment.

% [TODO] ADD CLOSING CLINICAL SIGNIFICANCE SENTENCE. Example:
% "Ultimately, a validated continuous monitoring system of this kind
% could reduce the burden on specialist clinics and enable earlier,
% more objective intervention for millions of patients living with
% Parkinson's Disease."


% ════════════════════════════════════════════════════════════════════
% IMPLEMENTATION DETAILS
%   [✓] Framework (PyTorch, SciPy, Optuna, Raspberry Pi)
%   [✓] Base architecture hyperparameters
%   [✓] Optuna-optimised architecture hyperparameters
%   [TODO] Move hardware spec (which GPU) — currently says "CUDA-enabled GPU v."
%          which is an incomplete sentence. Specify: e.g., NVIDIA RTX 3090.
%   [TODO] Consider moving this into Section V as \subsection{Experimental Setup}
%          rather than a standalone unnumbered section at the end.
% ════════════════════════════════════════════════════════════════════
\section*{Implementation Details}

All models are implemented in PyTorch and trained on a CUDA-enabled GPU. Signal processing uses SciPy (Butterworth filter, polyphase resampling). Hyperparameter search uses Optuna with the TPE sampler over 40 trials in a nested cross-validation framework, with Fold~1 as the internal validation fold. Edge inference runs on a Raspberry Pi CPU. Gradient clipping is applied at unit norm; the learning rate scheduler is ReduceLROnPlateau (factor $0.5$, patience $5$ epochs).

The base architecture uses $d_{\text{model}}{=}128$, $N{=}4$ cross-attention layers, $h{=}8$ attention heads, feed-forward dimension $512$, dropout $0.1$, and sequence length $T{=}256$. The edge-optimised variant selected by Optuna uses $d_{\text{model}}{=}32$, $N{=}3$, feed-forward dimension $256$, learning rate $2.91{\times}10^{-4}$, weight decay $1.62{\times}10^{-4}$, and batch size $32$.

% ════════════════════════════════════════════════════════════════════
% REFERENCES
%   [TODO] ALL \cite{ref} placeholders must be replaced with real entries.
%          Minimum required citations:
%          - PADS dataset (PhysioNet)
%          - Vaswani et al. 2017 (Transformer) — already labelled vaswani2017
%          - infoNCE: Oord et al. 2018 (CPC paper)
%          - TS-TCC: Eldele et al. 2023
%          - Time warping augmentation: Um et al. 2017
%          - LoRA: Hu et al. 2022
%          - TimesFM: Das et al. 2024 (Google)
%          - Nested CV: Wilimitis & Walsh 2023 or Varma & Simon 2006
%          - PD prevalence statistic (Introduction, first sentence)
%          - At least 2 prior IMU+PD papers for Related Work
% ════════════════════════════════════════════════════════════════════

% \bibliographystyle{IEEEtran}
% \bibliography{references}

\end{document}